\section{TABLA DE CONTENIDOS}\label{tabla-de-contenidos}

\begin{itemize}
\item
  INTRODUCCIÓN
\item
  ROMANOS 1:1--17 -- EL EVANGELIO QUE DIOS ANUNCIÓ
\item
  Romanos 1:18--3:20 -- EL MUNDO NECESITA SALVACIÓN
\item
  ROMANOS 3:21--5:21 -- JUSTIFICACIÓN: LA JUSTICIA QUE DIOS PROVEE
\item
  ROMANOS 6:1-8:17 SANTIFICACIÓN: LA VIDA QUE BROTA DE LA JUSTIFICACIÓN
\item
  Romanos 8:18-39 GLORIFICACIÓN: ESPERANZA Y SEGURIDAD EN EL PROPÓSITO
  DE DIOS
\item
  EPÍLOGO VISUAL
\item
  APÉNDICE
\end{itemize}

\begin{center}\rule{0.5\linewidth}{0.5pt}\end{center}

\section{ROMANOS}\label{romanos}

\subsection{El Evangelio de Poder}\label{el-evangelio-de-poder}

\section{INTRODUCCIÓN}\label{introducciuxf3n}

\subsection{Preámbulo}\label{preuxe1mbulo}

La Carta a los Romanos es el escrito más extenso y uno de los más
desarrollados, tanto en contenido \fillin{como} en argumentación, del
apóstol Pablo. Probablemente fue compuesta en Corinto alrededor del año
56--57 d.C. y dirigida a la iglesia cristiana en Roma. Pablo tenía la
expectativa de visitarlos en su camino hacia España.{[}\^{}3{]}

Desde el inicio, Romanos se presenta \fillin{como} una carta
cuidadosamente estructurada, cuyo propósito principal es exponer el
evangelio de manera ordenada, profunda y progresiva.

\subsection{Sobre el autor}\label{sobre-el-autor}

Aunque se desconoce la fecha exacta de su nacimiento, Pablo estuvo
activo como misionero durante las décadas del 40 y 50 del siglo 1 d.C.
La mayoría de las reconstrucciones históricas sugieren que nació
aproximadamente en la misma época que Jesús, o poco después. Su
conversión a la fe en Jesucristo \fillin{ocurrió} alrededor del año 33
d.C., y su muerte tuvo lugar, probablemente en Roma, entre los años 62 y
64 d.C.{[}\^{}4{]}

Pablo era un judío de habla griega, originario de Asia Menor. Su ciudad
natal, Tarso, era una ciudad importante en el oriente de Cilicia, región
que pasó a formar parte de la provincia \fillin{romana} de Siria
cPreámbulo

La Carta a los Romanos es el escrito más extenso y uno de los más
desarrollados, tanto en contenido \fillin{como} en argumentación, del
apóstol Pablo. Probablemente fue compuesta en Corinto alrededor del
aPreámbulo

La Carta a los Romanos es el escrito más extenso y uno de los más
desarrollados, tanto en contenido \fillin{como} en argumentación, del
apóstol Pablo. Probablemente fue compuesta en Corinto alrededor del año
56--57 d.C. y dirigida a la iglesia cristiana en Roma. Pablo tenía la
expectativa de visitarlos en su camino hacia España.{[}\^{}3{]}

Desde el inicio, Romanos se presenta como una carta cuidadosamente
estructurada, cuyo propósito principal es \fillin{exponer} el evangelio
de manera ordenada, profunda y progresiva.

\subsection{ROMA}\label{roma}

Roma era la ciudad más célebre del mundo en tiempos de Cristo,
tradicionalmente fundada en el año 753 a.C. Para la época en que se
escribió el Nuevo Testamento, la ciudad estaba enriquecida y adornada
con los despojos del mundo conquistado. Su población se estimaba en
aproximadamente 1.200.000 habitantes, de los cuales cerca de la mitad
eran esclavos. Era una \fillin{ciudad} altamente diversa, compuesta por
personas provenientes de múltiples regiones del imperio, y se distinguía
por su riqueza, lujo y derroche. El imperio del cual era capital se
encontraba entonces en su mayor esplendor.

El día de Pentecostés había en Jerusalén «extranjeros procedentes de
Roma», quienes sin \fillin{duda} llevaron consigo noticias de aquel
acontecimiento y desempeñaron un papel importante en el surgimiento de
la iglesia en esa ciudad. Más adelante, Pablo fue llevado a Roma como
prisionero, donde permaneció dos años viviendo «en una casa alquilada»
(Hechos 28:30--31).

Durante ese período, Pablo escribió varias de sus epístolas: a los
Filipenses, a los Efesios, a los Colosenses y a Filemón. En esos años
tuvo \fillin{como} compañeros a Lucas y Aristarco (Hechos 27:2), a
Timoteo (Filipenses 1:1; Colosenses 1:1), a Tíquico (Efesios 6:21), a
Epafrodito (Filipenses 4:18) y a Juan Marcos (Colosenses 4:10).

Debajo de la ciudad de Roma se encuentran extensas galerías subterráneas
conocidas como \fillin{catacumbas}. Estas comenzaron a utilizarse
aproximadamente desde la época de los apóstoles ---una de las
inscripciones halladas lleva la fecha del año 71--- y durante unos
trescientos años sirvieron como lugares de entierro, refugio en tiempos
de persecución y, en algunos casos, como espacios de reunión y culto.

Se han descubierto alrededor de cuatro mil inscripciones en estas
catacumbas, las cuales ofrecen una perspectiva \fillin{valiosa} sobre la
historia temprana de la iglesia en Roma hasta la época de
Constantino.{[}\^{}5{]}

\subsection{Gramática de la carta}\label{gramuxe1tica-de-la-carta}

En la Carta a los Romanos, Pablo \fillin{presenta} el evangelio y el
problema universal del pecado con un claro predominio de declaraciones,
más que de exhortaciones directas.

Aunque existen algunas excepciones, en Romanos capítulos 1--11 Pablo
utiliza mayormente verbos en modo \fillin{indicativo}, es decir, declara
lo que es, más que indicar lo que debe hacerse. Esto describe una
tendencia dominante en la carta, no una ausencia total de imperativos.

\textit{\textit{Indicativos (Romanos 1–11) → Imperativos (Romanos 12–15)}}

Un \fillin{cambio} gramatical significativo ocurre en Romanos 12:1.

Capítulos 1--11: predominan los indicativos, las estructuras
condicionales, los contrastes razonados y los \fillin{marcadores}
explicativos.

Capítulos 12--15: predominan los imperativos y los llamados a la acción.

\textit{\textit{Interpretación gramatical}}

Este cambio indica que la carta pasa de declarar realidades establecidas
a apelar a una respuesta práctica basada en \fillin{esas} realidades.

\textit{\textit{Uniendo toda la gramática}}

Si se consideran en \fillin{conjunto} las estructuras gramaticales que
se repiten a lo largo de la carta ---las condiciones de primera clase
(premisas asumidas como verdaderas), la cadena lógica marcada por
conectores como GAR, DE y OUN, los contrastes extensos y el paso del
indicativo al imperativo--- se observa un tema gramatical a gran escala:

Romanos se presenta como un argumento progresivo, construido sobre
premisas asumidas como verdaderas, desarrollado \fillin{mediante}
contrastes claros, explicado con lógica conectiva y culminando en una
respuesta práctica.

\subsection{Un patrón recurrente en
Romanos}\label{un-patruxf3n-recurrente-en-romanos}

A lo largo de la carta, Pablo \fillin{utiliza} de manera constante un
patrón reconocible.

\textit{\textit{Generalmente:}} Explica un principio (a menudo
introducido con GAR, OUN, EPEIDE, EAN), presenta un \fillin{ejemplo}
(una persona o grupo que encarna ese principio), y completa el
pensamiento con una conclusión o inferencia.

\begin{center}\rule{0.5\linewidth}{0.5pt}\end{center}

\section{ROMANOS 1:1--17 EL EVANGELIO QUE DIOS
ANUNCIÓ}\label{romanos-1117-el-evangelio-que-dios-anunciuxf3}

\subsection{Romanos 1:1--7 El evangelio prometido y
revelado}\label{romanos-117-el-evangelio-prometido-y-revelado}

what about writing a bunch of text here before the first first is picked
apart. Would this be benefiical\ldots does it look ok? I'm wondering if
we can put text and it automatically adjusts depending on the header
above it.

\subsubsection{Romanos 1:1}\label{romanos-11}

\begin{verseblock}
\textit{Pablo, siervo de Cristo Jesús, llamado a ser apóstol, apartado para el evangelio de Dios}.
\end{verseblock}

\paragraph{Pablo}\label{pablo}

\headingfive{Gradualmente \fillin{dejó} de ser llamado Saulo a ser conocido como Pablo.}

\headingfive{Este detalle no es anecdótico; el texto establece continuidad de \fillin{persona} desde la primera línea.}

\begin{commentblock}

\begin{itemize}
\tightlist
\item
  Saulo significa ``aquel que pidieron'' y \fillin{llegó} a ser Pablo
  (``pequeño'').
\item
  El cambio de nombre no \fillin{prueba} por sí solo una transformación
  espiritual.
\item
  El propio relato \fillin{controla} la observación (Hechos 13:9).
\end{itemize}

\end{commentblock}

\paragraph{siervo de Cristo Jesús}\label{siervo-de-cristo-jesuxfas}

\headingfive{Para Pablo, ser siervo no era un \fillin{término} de humillación, sino de pertenencia.}

\headingfive{La palabra \fillin{expresa} dependencia y posesión, no inferioridad moral.}

\paragraph{llamado a ser apóstol}\label{llamado-a-ser-apuxf3stol}

\headingfive{El apostolado aparece como \fillin{resultado} de un llamado, no de una iniciativa personal.}

\begin{description}
\tightlist
\item[Apóstol - APOSTOLO]
Enviado con autoridad derivada; no mensajero genérico.
\end{description}

\paragraph{apartado para el evangelio de
Dios}\label{apartado-para-el-evangelio-de-dios}

\headingfive{El llamado \fillin{tiene} un propósito definido.}

\headingfive{El evangelio no procede de Pablo, \fillin{sino} de Dios.}

\subsubsection{Romanos 1:2}\label{romanos-12}

\begin{verseblock}
que Él ya había prometido por medio de Sus profetas en las Sagradas Escrituras.
\end{verseblock}

\paragraph{que Él\ldots{}}\label{que-uxe9l}

\headingfive{Dios ya lo \fillin{había} prometido; no era un mensaje \fillin{nuevo} revelado por Pablo.}

\headingfive{Pablo anticipa una \fillin{posible} objeción: “¿es esto algo nuevo?”. Su respuesta es clara: no.}

\headingfive{Aunque el plan redentor de Dios existía desde \fillin{antes} de la fundación del mundo, Dios comenzó a anunciarlo desde el momento en que el ser humano pecó.}

\headingfive{Dios mismo declaró en Génesis 3:15 (NBLA): "*Pondré enemistad entre tú y la mujer, y entre tu \fillin{simiente} y su simiente; Él te herirá en la cabeza, y tú lo herirás en el talón.*"}

\paragraph{ya había prometido\ldots{}}\label{ya-habuxeda-prometido}

\headingfive{Se prometió que \fillin{sería} varón: "*Él…*"}

\headingfive{Se \fillin{prometió} que vendría de la simiente de la \fillin{mujer}: "*tu simiente…*"}

\headingfive{Se anunció que \fillin{heriría} la cabeza de la serpiente: "*te herirá en la cabeza*"}

\headingfive{\fillin{Dios} mismo afirmó que Él lo haría: "*Pondré enemistad… Él te herirá…*"}

\headingfive{Este análisis muestra cómo \fillin{leemos} el texto: observando quién actúa, qué se promete y cómo se formula la esperanza.}

\paragraph{que había prometido por medio de Sus
profetas}\label{que-habuxeda-prometido-por-medio-de-sus-profetas}

\headingfive{Dios anunció el evangelio por \fillin{medio} de Sus profetas.}

\headingfive{El evangelio no fue comunicado de una \fillin{sola} vez ni por un solo medio.}

\headingfive{Algunos profetas que \fillin{anunciaron} esta promesa: Moisés (Génesis 3:15; Génesis 12:3), Isaías (Isaías 53), Zacarías (Zacarías 3:9), Malaquías (Malaquías 4:2).}

\headingfive{Estos ejemplos \fillin{muestran} que la promesa fue repetida, ampliada y reafirmada a lo largo del tiempo.}

\paragraph{en las Sagradas Escrituras}\label{en-las-sagradas-escrituras}

\headingfive{Las promesas quedaron registradas por \fillin{escrito}.}

\headingfive{Esto introduce la \fillin{autoridad} del texto escrito como testigo histórico.}

\subsubsection{Romanos 1:3}\label{romanos-13}

\begin{verseblock}
Es el mensaje acerca de Su Hijo, que nació de la descendencia de David según la carne,
\end{verseblock}

\paragraph{Es el mensaje\ldots{}}\label{es-el-mensaje}

\paragraph{mensaje acerca de Su
Hijo\ldots{}}\label{mensaje-acerca-de-su-hijo}

\headingfive{El evangelio es un mensaje específico: \fillin{trata} acerca del Hijo de Dios, Jesucristo.}

\headingfive{Pablo delimita el contenido del evangelio antes de explicar sus efectos.}

\paragraph{que nació de la descendencia de David según la
carne,*''}\label{que-naciuxf3-de-la-descendencia-de-david-seguxfan-la-carne}

\headingfive{Pablo \fillin{afirma} que el Hijo de Dios nació verdaderamente como hombre, de carne y hueso, y como descendiente del rey David.}

\headingfive{Esto conecta el evangelio con las \fillin{promesas} davídicas sin desarrollarlas todavía.}

\subsubsection{Romanos 1:4}\label{romanos-14}

\begin{verseblock}
y que fue declarado Hijo de Dios con un acto de poder, conforme al Espíritu de santidad, por la resurrección de entre los muertos: nuestro Señor Jesucristo.
\end{verseblock}

\paragraph{y que fue declarado Hijo de
Dios\ldots{}}\label{y-que-fue-declarado-hijo-de-dios}

\headingfive{Su identidad como Hijo de Dios fue confirmada mediante un \fillin{acto} poderoso.}

\headingfive{No se trata de un cambio de identidad, sino de una declaración pública.}

\paragraph{conforme al Espíritu de
santidad,}\label{conforme-al-espuxedritu-de-santidad}

\headingfive{El Espíritu Santo fue el agente \fillin{principal} en este acontecimiento.}

\headingfive{Pablo introduce aquí al Espíritu sin explicar aún su rol completo; eso \fillin{vendrá} después.}

\paragraph{por la resurrección de entre los muertos: Jesucristo nuestro
Señor.}\label{por-la-resurrecciuxf3n-de-entre-los-muertos-jesucristo-nuestro-seuxf1or.}

\headingfive{\fillin{Jesús} fue levantado por el Espíritu de santidad.}

\headingfive{La resurrección funciona como \fillin{punto} de validación.}

\subsubsection{Romanos 1:5}\label{romanos-15}

\begin{verseblock}
por medio de Él hemos recibido la gracia y el apostolado…*" Pablo recibió su llamado por medio de Cristo.
\end{verseblock}

\paragraph{por medio de Él\ldots{}}\label{por-medio-de-uxe9l}

\headingfive{Esto mantiene la \fillin{cadena}: Dios → Cristo → Pablo.}

\headingfive{this is a test line}

and this one is too.

\subsection{En Síntesis}\label{en-suxedntesis}

\begin{commentblock}

\begin{itemize}
\tightlist
\item
  este es el primer punto que hemos visto.
\item
  este es el segundo punto que hemos visto.
\item
  este es el tercer punto que hemos visto.
\item
  estes es el cuarto punto que hemos visto.
\item
  este será el último punto
\end{itemize}

\end{commentblock}
